
\section{4. Method}

\subsection{4.1 Overview}

NHA runs two attention streams in parallel and combines them via a token-level 
gate:

1. \textbf{Local stream} $O^\text{local}$: sliding-window attention.
2. \textbf{Global stream} $O^\text{global}$: RFF linear attention.
3. \textbf{Gate} $g$: per-token sigmoid.
4. \textbf{Fusion} $O = g \odot O^\text{local} + (1 - g) \odot O^\text{global}$.

\textbf{Figure 4.1} illustrates the pipeline.

\begin{lstlisting}
        Q, K, V $\in$ R^{B $\times$ H $\times$ L $\times$ d_h}
                 |
        +--------+--------+
                           
   +----------+      +----------+
   |  Local   |      |  Global  |
   |  window  |      |   RFF    |
   | O_local  |      | O_global |
   +----+-----+      +----+-----+
        |                 |
        +--------+--------+
                  
        +-----------------+
        |  Gate g(x)      |
        |  g O_l+(1-g) O_g|
        +--------+--------+
                  
              out_proj
\end{lstlisting}

\subsection{4.2 Local Window Stream}

\subsubsection{4.2.1 Definition}

For window size $w$:

\[
S_{ij} = \frac{Q_i^\top K_j}{\sqrt{d_h}}, \quad
\text{mask}_{ij} = \mathbb{1}[|i - j| \leq w]
\tag{4.1}
\]

\[
O_i^\text{local} = \text{softmax}(S \odot \text{mask})_i \cdot V
\tag{4.2}
\]

\subsubsection{4.2.2 Chunked Computation}

To bound memory, we process the query dimension in chunks of 64:

\begin{lstlisting}
for i in range(0, N, chunk):
    end = min(i + chunk, N)
    q_block = Q[:, :, i:end]
    k_start = max(0, i - window)
    k_end = min(N, end + window)
    k_block = K[:, :, k_start:k_end]
    # ... compute scores, mask, softmax ...
\end{lstlisting}

This gives $\mathcal{O}(B H C (2w + C))$ peak memory instead of 
$\mathcal{O}(B H L^2)$.

\subsubsection{4.2.3 Causal Masking}

For decoders:

\[
\text{mask}_{ij}^\text{causal} = \mathbb{1}[j \leq i \text{ and } i - j \leq w]
\tag{4.3}
\]

\subsection{4.3 Global RFF Stream}

\subsubsection{4.3.1 RFF Projection}

We build $\Omega \in \mathbb{R}^{d_h \times m}$ via block QR:

\begin{lstlisting}
blocks = []
remaining = m
while remaining > 0:
    block_size = min(d_h, remaining)
    g = torch.randn(d_h, block_size)
    q, _ = torch.linalg.qr(g)
    blocks.append(q)
    remaining -= block_size
W = torch.cat(blocks, dim=1) * math.sqrt(d_h)
\end{lstlisting}

\textbf{Orthogonality.} Each block $Q_i$ has orthonormal columns. The scale 
$\sqrt{d_h}$ matches the expected variance of Gaussian projections.

\subsubsection{4.3.2 Feature Map}

\[
\phi(x) = \frac{1}{\sqrt{m}} \left(\mathrm{elu}(\Omega^\top x) + 1\right)
\tag{4.4}
\]

\textbf{Properties:}
\begin{itemize}
\item \textbf{Non-negative:} $\mathrm{elu}(x) > -1$ $\to$ $\phi > 0$.
\end{itemize}
\begin{itemize}
\item \textbf{Bounded:} $\phi \in [0, \infty)$.
\end{itemize}
\begin{itemize}
\item \textbf{Smooth:} ELU is $C^1$.
\end{itemize}

\subsubsection{4.3.3 Non-Causal Computation}

For encoders:

\[
KV = \phi(K)^\top V \in \mathbb{R}^{m \times d_h}, \quad
Z = \sum_j \phi(K_j) \in \mathbb{R}^m
\tag{4.5}
\]

\[
O_i^\text{global} = \frac{\phi(Q_i)^\top KV}{\phi(Q_i)^\top Z + \epsilon}
\tag{4.6}
\]

\subsubsection{4.3.4 Causal Prefix Scan (Chunked)}

For decoders, we use a chunked prefix scan with cumulative state:

\[
S \in \mathbb{R}^{B \times H \times m \times d_h}, \quad
Z \in \mathbb{R}^{B \times H \times m}
\]

Initialize $S = 0$, $Z = 0$. For each chunk $[start, end)$:

1. Compute intra-chunk causal scores:
\[
   A^\text{intra} = \phi(Q_c) \phi(K_c)^\top \odot \text{tril}(1)
   \tag{4.7}
\]

2. Compute inter-chunk contribution:
\[
   O^\text{inter} = \phi(Q_c) S
   \tag{4.8}
\]

3. Combine:
\[
   O_c = \frac{A^\text{intra} V_c + O^\text{inter}}{(A^\text{intra} \cdot 1 + \phi(Q_c) Z) + \epsilon}
   \tag{4.9}
\]

4. Update state:
\[
   S \leftarrow S + \phi(K_c)^\top V_c, \quad Z \leftarrow Z + \sum \phi(K_c)
   \tag{4.10}
\]

\textbf{Complexity.} $\mathcal{O}(B H L m d_h)$ time, $\mathcal{O}(B H m d_h)$ memory.

\subsubsection{4.3.5 Numerical Stability}

We compute the RFF stream in \textbf{FP32} inside an autocast-disabled block:

\begin{lstlisting}
with torch.autocast(device_type=..., enabled=False):
    phi_Q_f = phi_Q.float()
    phi_K_f = phi_K.float()
    V_f = V_masked.float()
    # ... computation ...
\end{lstlisting}

This avoids FP16 overflow for $L > 512$.

\subsection{4.4 Token-Level Gate}

\subsubsection{4.4.1 Parameterization}

\[
g_i = \sigma(w^\top x_i + b)
\tag{4.11}
\]

where $x_i$ is the input hidden state at position $i$, and $w \in \mathbb{R}^d$, 
$b \in \mathbb{R}$.

\subsubsection{4.4.2 Initialization}

We initialize $w = 0$ and $b = 3.0$, so:

\[
g_i = \sigma(0 + 3) = \sigma(3) \approx 0.95 \quad \text{for all } i
\tag{4.12}
\]

\textbf{Why bias toward local?} The local window matches the pretrained model's 
short-range behavior. Starting global-leaning would immediately degrade 
accuracy.

\subsubsection{4.4.3 Per-Token vs Per-Head}

NHA uses \textbf{per-token} gating (single scalar per position), while HALO-ReGLA 
uses per-head. Per-token is:

\begin{itemize}
\item \textbf{Parameter-efficient:} $d + 1$ parameters instead of $d H + H$.
\end{itemize}
\begin{itemize}
\item \textbf{Position-aware:} Can adapt to local vs global needs per position.
\end{itemize}
\begin{itemize}
\item \textbf{Simpler:} No head dimension to manage.
\end{itemize}

Empirically, per-token works as well as per-head for our tasks.

\subsection{4.5 Complete Forward Pass}

\textbf{Algorithm 1:} NHA forward
\begin{lstlisting}
Input:  hidden_states $\in$ R^{B $\times$ L $\times$ d}
        attention_mask, position_embeddings, is_causal
        window_size w, rff_features m
Output: output $\in$ R^{B $\times$ L $\times$ d}

 1: Q, K, V $\leftarrow$ proj(hidden_states)
 2: Q, K $\leftarrow$ apply_rope(Q, K, position_embeddings)
 3: K, V $\leftarrow$ repeat_kv(K, V, H / H_kv) if GQA
 4:
 5: # Local stream (chunked)
 6: O_local $\leftarrow$ sliding_window_attention(Q, K, V, w, is_causal)
 7:
 8: # Global RFF stream
 9:  _Q $\leftarrow$ elu($\Omega$^T Q) + 1
10:  _K $\leftarrow$ elu($\Omega$^T K) + 1
11: if is_causal:
12:     (S, Z) $\leftarrow$ initialize_cumulative_state()
13:     for chunk in chunks(L, 64):
14:         O_global[chunk] $\leftarrow$ RFF_chunk( _Q,  _K, V, S, Z)
15:         (S, Z) $\leftarrow$ update_state(S, Z,  _K[chunk], V[chunk])
16: else:
17:     KV $\leftarrow$  _K^T V
18:     Z $\leftarrow$ sum( _K)
19:     O_global $\leftarrow$  _Q @ KV / ( _Q @ Z + $\epsilon$)
20:
21: # Token-level gate
22: g $\leftarrow$ sigmoid(gate(hidden_states))         # (B, L, 1)
23: O $\leftarrow$ g * O_local + (1 - g) * O_global
24: O $\leftarrow$ O * valid_mask
25: O $\leftarrow$ reshape and out_proj
26: return O
\end{lstlisting}

\subsection{4.6 Complexity Analysis}

\begin{table}[H]
\centering
\begin{tabular}{llll}
\toprule
\textbf{Component} & \textbf{Local} & \textbf{Global RFF} & \textbf{NHA (total)} \\
\midrule
Time & $\mathcal{O}(L w d_h)$ & $\mathcal{O}(L m d_h)$ & $\mathcal{O}(L (w + m) d_h)$ \\
Memory & $\mathcal{O}(L w)$ & $\mathcal{O}(m d_h)$ & $\mathcal{O}(L w + m d_h)$ \\
vs Dense & $\mathcal{O}(L^2 d_h)$ & -- & **$\mathcal{O}(L (w + m) d_h)$** \\
\bottomrule
\end{tabular}
\end{table}

\textbf{Key insight.} NHA is \textbf{linear} in $L$ for both branches. This is the 
primary efficiency benefit.

\subsection{4.7 Parameter Analysis}

\begin{table}[H]
\centering
\begin{tabular}{lll}
\toprule
\textbf{Component} & \textbf{Baseline} & \textbf{NHA} \\
\midrule
$W_Q$ & $d H d_h$ & Same \\
$W_K$ & $d H_{kv} d_h$ & Same \\
$W_V$ & $d H_{kv} d_h$ & Same \\
$W_O$ & $H d_h d$ & Same \\
**Gate** & 0 & $d + 1$ \\
**RFF proj** & 0 & $d_h m$ (buffer, not trainable) \\
LoRA & 0 & $2 r (\cdots)$ \\
\bottomrule
\end{tabular}
\end{table}

\textbf{Total: same as MHA + $d + 1$ parameters (gate) + RFF buffer.}

\subsection{4.8 Theoretical Notes}

\subsubsection{4.8.1 Kernel Approximation}

With orthogonal $\Omega$ and $m \to \infty$:

\[
\phi(q)^\top \phi(k) \xrightarrow{m \to \infty} \mathbb{E}_\omega[\ldots] = k_\text{exp}(q, k)
\tag{4.13}
\]

where $k_\text{exp}$ is the exp kernel.

\subsubsection{4.8.2 Gate Convergence}

The gate converges to a stationary point where:

\[
\frac{\partial \mathcal{L}}{\partial g_i} = (\nabla_O \mathcal{L}) \cdot (O^\text{local}_i - O^\text{global}_i) = 0
\tag{4.14}
\]

or hits the boundary ($g_i \to 0$ or $g_i \to 1$).

\subsubsection{4.8.3 Causality}

For causal models, the chunked prefix scan ensures:

\[
O_i^\text{global} = f(\phi(Q_j), \phi(K_j), V_j : j \leq i)
\tag{4.15}
\]

No future information leaks.